% Proof corresponding exactly to statement.tex in this isolated dossier.
% No production call has been made for this follow-up candidate.

\begin{proof}
We use the following already frozen and production-verified local result,
denoted by \((\mathrm M)\).  If \(Z\) is locally compact Hausdorff with a
clopen basis and \(P\) is any nonempty convex polyhedron, including an
unbounded, lower-dimensional, or zero-dimensional one, then on
\(Z\times P\) the first local decreasing-limit stage of the
single-chamber smooth psh generators, the full decreasing-sequentially
generated psh sheaf, and the Stone-factorized fibrewise-convex comparison
sheaf coincide.  Moreover, every germ has a decreasing approximation by
smooth psh generators on one fixed product
\(C\times W\), where \(C\) is compact clopen and \(W\) is a relatively
open convex polyhedral neighbourhood.  The exact statement, proof, source
chain, hostile audit, and production receipt for \((\mathrm M)\) are pinned
in \texttt{SOURCE\_AND\_OBLIGATION\_MATRIX.md}.

We first record the elementary transport argument needed to apply
\((\mathrm M)\).  Let \(h:T\to T'\) be a homeomorphism, and let
\(\mathcal G_0\) and \(\mathcal G'_0\) be generator sheaves such that
\begin{equation}
                   f\in\mathcal G_0
             \quad\Longleftrightarrow\quad
                   f\circ h^{-1}\in\mathcal G'_0
 \label{eq:p071-complex-stage-zero-transport}
\end{equation}
locally, in both directions.  A local sequence
\(f_j\in\mathcal G_0\) with \(f_j\downarrow f\) transports term by term to
\(f_j\circ h^{-1}\downarrow f\circ h^{-1}\).  Applying the same argument
to \(h^{-1}\) proves that \(h\) identifies the two first local
decreasing-limit stages.

It also identifies the full generated sheaves.  Indeed, the inverse image
under \(h\) of the full generated sheaf on \(T'\) is a sheaf containing
\(\mathcal G_0\) and stable under local decreasing sequential limits.
Minimality gives one inclusion; applying this to \(h^{-1}\) gives the
reverse inclusion.  Thus
\begin{equation}
 h_*(\mathcal G_1)=\mathcal G'_1,
 \qquad
 h_*(\mathcal G)=\mathcal G'.
 \label{eq:p071-complex-all-stage-transport}
\end{equation}
This uses neither a partition of unity nor a countability property of the
underlying spaces.

We will also use that this construction commutes with restriction to an
open subset.  The assertion is immediate for the first local stage.  For
the full sheaf, let \(S\subset T\) be open and write \(\mathcal G_S\) for
the sheaf generated on \(S\) by \(\mathcal G_0|_S\).  One inclusion follows
because \(\mathcal G|_S\) is decreasing-sequentially closed and contains
the restricted generators.  For the reverse inclusion, consider the
subsheaf of \(\mathcal G\) consisting of sections whose germs at points of
\(S\) lie in \(\mathcal G_S\), with no extra condition away from \(S\).
This subsheaf contains \(\mathcal G_0\) and is stable under local decreasing
sequential limits.  The minimality of \(\mathcal G\) therefore forces every
germ of \(\mathcal G|_S\) to lie in \(\mathcal G_S\).  Hence
\begin{equation}
                    \mathcal G|_S=\mathcal G_S.
 \label{eq:p071-complex-restriction}
\end{equation}

Fix \((z_0,p_0)\in X\), and choose a chart
\(\varphi:U\to V\subset P\) from
\eqref{eq:p071-complex-chart}.  On a clopen neighbourhood in \(Z\), the map
\begin{equation}
               \operatorname{id}\times\varphi:
               Z\times U\longrightarrow Z\times V
 \label{eq:p071-complex-product-chart}
\end{equation}
identifies the stage-zero sheaves by hypothesis.  Consequently
\eqref{eq:p071-complex-all-stage-transport} identifies both the first and
the full psh sheaves with their restrictions from the product model
\(Z\times P\), using \eqref{eq:p071-complex-restriction} for the open chart
target \(Z\times V\).

We next check that the comparison condition does not depend on the chart.
Suppose first that a germ has the factorization
\eqref{eq:p071-complex-comparison} in a chart \(\varphi_\alpha\).  By
\((\mathrm M)\), its coordinate representative belongs to the first-stage
model sheaf.  On an overlap, the transition map is generator-preserving;
the first equality in \eqref{eq:p071-complex-all-stage-transport} therefore
places the representative in the first-stage model sheaf of any other
chart \(\varphi_\beta\).  A second application of \((\mathrm M)\), after a
convex shrinking around the image of the point, produces exactly a compact
metrizable zero-dimensional quotient and an upper-semicontinuous
fibrewise-convex factorization in the \(\varphi_\beta\)-coordinates.  The
same argument with \(\alpha\) and \(\beta\) interchanged proves equivalence.
The comparison property is therefore intrinsic.  Since its condition is
local in any chart, it defines a sheaf \(\mathscr C_X\).

Now let \(u\) be a germ of \(\operatorname{Psh}_X\) at
\((z_0,p_0)\).  Transport through
\eqref{eq:p071-complex-product-chart} and
\eqref{eq:p071-complex-all-stage-transport}.  The model theorem
\((\mathrm M)\) places the transported germ in the model comparison sheaf,
and hence in the model first-stage sheaf.  Transporting back proves
\begin{equation}
                  \operatorname{Psh}_X
                  \subseteq\mathscr C_X
                  \subseteq\operatorname{Psh}_{X,1}.
 \label{eq:p071-complex-two-inclusions}
\end{equation}
The tautological inclusion
\(\operatorname{Psh}_{X,1}\subseteq\operatorname{Psh}_X\) completes the
proof of \eqref{eq:p071-complex-collapse}.

The fixed-neighbourhood assertion is equally local but is worth spelling
out.  Apply the fixed-domain clause of \((\mathrm M)\) to the coordinate
representative of \(u\).  After shrinking inside \(V\), it gives a compact
clopen \(C\), a relatively open convex polyhedral neighbourhood \(W\), and
model generators \(g_k\) on the one set \(C\times W\), with
\(g_{k+1}\leq g_k\) and \(g_k\downarrow u\) there.  Set
\begin{equation}
        G_k(z,p)=g_k\bigl(z,\varphi(p)\bigr).
 \label{eq:p071-complex-pulled-approximants}
\end{equation}
The generator-preserving condition makes every \(G_k\) a stage-zero
generator on the same set
\(C\times\varphi^{-1}(W)\), and order and pointwise convergence are
preserved.  This proves \eqref{eq:p071-complex-fixed-star}.  Since
\(\varphi^{-1}(W)\) is one ordinary neighbourhood of \(p_0\), it contains
the portions of every cell incident to the germ.  Thus the model
construction smooths across all chart faces simultaneously.  Separate
chamberwise sequences are never chosen and hence do not have to be glued.

This also explains the loop case.  The proof chooses only one chart at the
given germ.  Continuing that chart around a loop is unnecessary, so the
composition of affine transition maps around the loop does not enter any
step.  Conversely, nothing in the argument makes the local sequences agree
term by term on distinct chart overlaps.  It therefore proves the sheaf
identity, whose first stage is local by definition, but it proves no
single-global-sequence assertion.

We verify the asserted obstruction to weakening the hypotheses.  The
two-ray complex
\[
        L=\mathbf R_{\geq0}e_1\cup\mathbf R_{\geq0}e_2
\]
is face-to-face, and its vertex is incident to exactly two one-dimensional
cells, so its underlying space is a topological one-manifold.  The Hessian
of the function \(F\) in \eqref{eq:p071-complex-v-function} is zero.
Hence its restriction is an ambient-induced cellwise smooth psh generator.
It is continuous and the two restrictions agree at the common face.

Any cellwise affine homeomorphism from a neighbourhood of the vertex to an
interval sends the two rays to opposite half-intervals.  After possibly
interchanging them, it has the form
\[
                 te_1\longmapsto at,
        \qquad  te_2\longmapsto-bt
        \qquad(a,b>0).
\]
In the interval coordinate \(s\), the transported restriction of \(F\)
has slope \(b^{-1}>0\) to the left of zero and slope \(-a^{-1}<0\) to the
right.  A convex function must have nondecreasing one-sided slopes, so this
function is not convex; for \(a=b=1\) it is \(-|s|\).  Therefore no such
unfolding identifies this complete cellwise generator germ with the
single-convex-interval generator germ.  Agreement only on the common
zero-dimensional face does not impose the missing slope inequality.

Finally, all dimensional edge cases are inherited without change from
\((\mathrm M)\).  When \(P\) has dimension zero, its affine span is a
point, fibrewise convexity is vacuous, and the model proof uses no maximum
over an empty set.  For a lower-dimensional \(P\), every Hessian and every
relative interior is taken in \(\operatorname{Aff}(P)\).  At a boundary or
corner, \(W\) is relatively open in \(P\), so the fixed domain remains an
ordinary neighbourhood in \(|\Sigma|\).  Unboundedness causes no problem
because \((\mathrm M)\) constructs the approximation on a bounded compact
core after a local shrinking.  None of these observations supplies a chart
at a branched or singular-link germ, which is why those cases remain outside
the assertion.
\end{proof}
